\section*{Capítulo \ref{ch:b2:limb-organogenesis} --- Organogénesis: la construcción de la extremidad de los tetrápodos}

\begin{solution}{exo:b2:limb-organogenesis:1}
Proximodistal: la \omterm{def:b2:limb-organogenesis:bud}{cresta ectodérmica apical} a lo largo del extremo, los
FGF. Anteroposterior: la \omterm{def:b2:limb-organogenesis:bud}{zona de actividad polarizante} en el margen
posterior, \omterm{prop:b2:limb-organogenesis:zpa}{Sonic hedgehog}. Dorsoventral: el ectodermo dorsal, Wnt7a.
\end{solution}

\begin{solution}{exo:b2:limb-organogenesis:2}
Estilopodio: húmero, fémur (Hoxa/d 9--10). Zeugopodio: radio y cúbito,
tibia y peroné (Hox 11). Autopodio: muñeca y mano, tobillo y pie
(Hox 13).
\end{solution}

\begin{solution}{exo:b2:limb-organogenesis:3}
Se especifica el campo de la extremidad (Tbx5); el esbozo crece bajo el
FGF de la cresta; la ZPA y el ectodermo dorsal lo organizan; las células
que abandonan la \omterm{prop:b2:limb-organogenesis:aer}{zona de progreso} se condensan en radios; el cartílago se
diferencia; las articulaciones se forman donde se suprime el cartílago;
los vasos sanguíneos invaden y el hueso sustituye al cartílago; las
células interdigitales mueren; los precursores musculares de los somitos
migran hacia la extremidad, se fusionan y se unen mediante tendones.
\end{solution}

\begin{solution}{exo:b2:limb-organogenesis:4}
(a) Un muñón con poco o ningún húmero; (b) húmero, radio y cúbito, pero
sin mano; (c) una extremidad completa --- el FGF sustituye a la cresta.
\end{solution}

\begin{solution}{exo:b2:limb-organogenesis:5}
$x = 40\ln 2 = \qty{28}{\micro m}$, $40\ln 5 = \qty{64}{\micro m}$,
$40\ln 20 = \qty{120}{\micro m}$: territorios de 28, 36 y
\qty{56}{\micro m}.
\end{solution}

\begin{solution}{exo:b2:limb-organogenesis:6}
(a) 4-3-2-2-3-4, un juego especular completo. (b) Una fuente más débil da
un máximo más bajo, de modo que solo se añade el dedo de umbral bajo:
2-2-3-4. (c) Una exposición breve también especifica solo el dedo
adicional más anterior: 2-2-3-4 --- las células integran la concentración
a lo largo del tiempo.
\end{solution}

\begin{solution}{exo:b2:limb-organogenesis:7}
$5000\times 2^{8} = 1.3\times 10^{6}$. Llegar a $4\times 10^{6}$ requiere
$\log_2 800 = 9.6$ duplicaciones, de modo que el ciclo debe durar en
promedio unas \qty{10}{h}, o bien el recuento debe incluir los
precursores musculares que migran desde los somitos; la muerte celular
agravaría la discrepancia en lugar de reducirla.
\end{solution}

\begin{solution}{exo:b2:limb-organogenesis:8}
Las células interdigitales del pollo reciben una señal BMP y mueren; las
del pato expresan un inhibidor de la BMP y sobreviven, de modo que la
membrana persiste. El inhibidor en una bola salva la membrana del pollo:
una pata palmeada. Que la muerte se produzca a su hora en un injerto
significa que las células ya habían sido instruidas antes de la
extracción --- la decisión se toma pronto y se ejecuta más tarde, de forma
autónoma.
\end{solution}

\begin{solution}{exo:b2:limb-organogenesis:9}
Sin Hoxa13 ni Hoxd13: sin autopodio --- la extremidad anterior termina en
la muñeca. Sin Hoxa11 ni Hoxd11: un radio y un cúbito muy acortados, con
una mano unida casi directamente al húmero. Cada segmento necesita su
propio par Hox; los dominios no son meros marcadores, sino instrucciones.
\end{solution}

\begin{solution}{exo:b2:limb-organogenesis:10}
El FGF de la cresta mantiene activo Shh en la ZPA; Shh, a través de un
relevo mesodérmico, mantiene activo el FGF en la cresta. La señal de cada
centro certifica, por tanto, que el otro funciona, de modo que una
extremidad nunca crece sin organizarse ni se organiza sin crecer, y una
pérdida transitoria de cualquiera de los dos la repara el otro. Si Shh
no pudiera mantener el FGF, la cresta regresaría en un día y el
crecimiento se detendría en el segmento alcanzado: una extremidad
truncada. El bucle del parto lo termina el acontecimiento que impulsa
--- el nacimiento elimina el estímulo; el bucle de la extremidad lo
termina la diferenciación --- el mesodermo deja de hacer de relevo cuando
la extremidad está completa.
\end{solution}

\begin{solution}{exo:b2:limb-organogenesis:11}
$e^{4r} = 8$: $r = \ln 8/4 = \qty{0.52}{d^{-1}}$. Crecimiento: una
proliferación más larga y más rápida del cartílago de los dedos (un
potenciador de la extremidad que aumenta la expresión de un gen de
crecimiento en los dedos). Muerte: la supervivencia del tejido
interdigital gracias a la expresión de un inhibidor de la BMP, que
conserva la membrana.
\end{solution}

\begin{solution}{exo:b2:limb-organogenesis:12}
Las señales se reutilizan porque un módulo receptor--transducción que
funciona es barato de desplegar de nuevo: lo que cambia de un tejido a
otro no es la señal, sino el conjunto de genes que sus células receptoras
son competentes para activar, de modo que el mismo perfil de
concentración significa motoneurona en un lugar y dedo 4 en otro. Una
\omterm{def:b2:genome-diversification:mutation}{mutación} en un gen así afecta por tanto a muchos órganos a la vez --- las
\omterm{def:b2:genome-diversification:mutation}{mutaciones} de Shh dan a la vez holoprosencefalia y defectos de las
extremidades ---, salvo que afecte a un potenciador específico de tejido,
en cuyo caso solo cambia un órgano, que es como se produce la mayor parte
del cambio evolutivo en estos genes.
\end{solution}

\begin{solution}{pb:b2:limb-organogenesis:1}
\textbf{1.} $96/12 = 8$ duplicaciones: $2000\times 256 = 5.1\times 10^{5}$
células.
\textbf{2.} $1000/256 \approx 4$: aumento de tamaño celular y matriz
cartilaginosa.
\textbf{3.} Longitud $\times 10$; un cilindro de radio constante daría
$\times 10$ en volumen; el factor cien adicional significa que el radio
también creció unas diez veces --- el esbozo se ensanchó en forma de pala
a medida que se alargaba.
\textbf{4.} A los 3 días, $300/400 = \qty{75}{\%}$; a los 7 días, $300/4000 =
\qty{7.5}{\%}$.
\textbf{5.} Solo el extremo está bajo la señal; a medida que el extremo
avanza, las células que quedan atrás abandonan la zona en el orden en que
se formaron, de modo que las estructuras proximales se especifican
primero y las distales al final.
\textbf{6.} El zeugopodio (radio y cúbito), especificado entre los 3,5 y
los 4 días.
\textbf{7.} Estilopodio hacia los 3,5 días, zeugopodio hacia los 4,0 días,
autopodio hacia los 4,5 días (las puntas de los dedos, más tarde).
\textbf{8.} Medio día: un ciclo celular.
\textbf{9.} Cada ciclo pasado en la \omterm{prop:b2:limb-organogenesis:aer}{zona de progreso} hace avanzar el
código Hox (9--10, después 11, después 13); las células que salen tras
uno, dos o tres ciclos llevan el código del estilopodio, del zeugopodio o
del autopodio.
\textbf{10.} Un ala completa: el FGF es lo que aportaba la cresta.
\textbf{11.} Un segundo crecimiento desde la superficie dorsal:
estructuras distales duplicadas, una extremidad supernumeraria.
\textbf{12.} El mesodermo deja de transmitir Shh a la cresta, la
expresión de Shh cesa, la cresta regresa y las células se diferencian en
lugar de dividirse: el bucle que sostenía el crecimiento se rompe.
\textbf{13.} $k = \ln 2/1740 = \qty{4.0e-4}{s^{-1}}$; $\lambda =
\sqrt{10^{-12}/4\times 10^{-4}} = \qty{50}{\micro m}$.
\textbf{14.} $x = 50\ln(1/0.6) = 26$, $50\ln 4 = 69$, $50\ln 12.5 =
\qty{126}{\micro m}$: territorios de 26, 43 y \qty{57}{\micro m}.
\textbf{15.} $600 - 126 = \qty{474}{\micro m}$.
\textbf{16.} Todas las fronteras se desplazan hacia fuera en
$50\ln 2 = \qty{35}{\micro
m}$: 61, 104, 161. Las tres anchuras son 61, 43 y \qty{57}{\micro m}:
solo crece el territorio del dedo 4, los demás simplemente se desplazan.
El patrón se corre hacia delante en lugar de ganar un dedo, y la región
anterior sin dedos se reduce de 474 a \qty{439}{\micro m}.
\textbf{17.} Cada juego necesita \qty{126}{\micro m}: al menos
\qty{252}{\micro m}; el esbozo de \qty{600}{\micro m} tiene sitio, con
\qty{348}{\micro m} de tejido sin dedos en el centro.
\textbf{18.} Con \qty{200}{\micro m} los dos gradientes se solapan y se
suman: las células centrales reciben suficiente para el dedo 3 y los
territorios del dedo 2 desaparecen --- 4-3-4 o 4-3-3-4, menos dedos que un
juego especular completo.
\textbf{19.} La cuarta parte de las células da $c_0/4$: ninguna célula
alcanza el umbral del dedo 4 ni el del 3 ($0.25 c_0$ en el propio margen),
y el umbral del dedo 2 se supera hasta $50\ln(0.25/0.08) =
\qty{57}{\micro m}$: un dedo 2 adicional.
\textbf{20.} $3\times 300\times 300\times 200 = 5.4\times 10^{7}$
\unit{\micro m^3}: $5.4\times 10^{4}$ células; a lo largo de
\qty{86400}{s}, $0.6$ células por segundo.
\textbf{21.} Bola de BMP en la membrana del pato: la membrana muere y los
dedos se separan; inhibidor en la membrana del pollo: la membrana
sobrevive, una pata palmeada.
\textbf{22.} $e^{0.5\times 4} = e^{2} = 7.4$. El segundo cambio es la
supervivencia de la membrana interdigital gracias a un inhibidor de la
BMP.
\textbf{23.} La fase tardía de expresión de Hox 13 que especifica un
autopodio, y la persistencia de la cresta: el pliegue del pez forma en su
lugar radios de aleta. La aleta de \emph{Tiktaalik} contiene huesos de
tipo carpiano --- una aleta que empieza a ser una extremidad.
\textbf{24.} $5.4\times 10^{4}/1000 = 54$: $\log_2 54 = 5.8$ duplicaciones,
unas \qty{69}{h}, tres días para formar lo que un día elimina.
\textbf{25.} Ocho duplicaciones; fronteras de los dedos en 26, 69 y
\qty{126}{\micro m}; una duplicación especular necesita al menos
\qty{252}{\micro m}.
\end{solution}
